\section{库存轨迹与服务水平验证}

\subsection{库存的业务含义}

库存变量$b_{dh}$不是仓库存货，而是截至小时$h$尚未完成分拣的货物。由式\eqref{eq:flow}，若某小时到货量突然增加，模型可以选择在后续班次中处理，但必须在第23小时将库存降为0。相比直接比较累计产能，库存轨迹能揭示拥堵发生的时刻和消化速度。

\begin{figure}[htbp]
\centering
\includegraphics[width=.88\textwidth]{figures/backlog_profile.pdf}
\caption{问题二平均库存轨迹与10\%--90\%分位带}
\label{fig:backlog}
\end{figure}

图\ref{fig:backlog}中16时虚线对应问题二的早间服务截止。程序对每一天单独检查式\eqref{eq:deadline}，因此平均图形只能作为总体展示，不能替代逐日核验；逐日结果已经在附录的720行明细表中给出。

\subsection{处理量与产能的差值}

定义小时余量
\begin{equation}
e_{dh}=\sum_s\delta_{sh}(25y_{ds}-15z_{ds,h-s})-p_{dh}\ge0.
\end{equation}
当$e_{dh}$较大时代表该小时存在冗余产能，当$e_{dh}=0$时代表产能成为瓶颈。余量与库存应联合解读：产能瓶颈不一定导致库存上升，因为此前可能没有货物到达；反之，库存上升说明此前累计可用处理量不足。

\subsection{截止约束的逐日核验}

对每一天计算早间到货量$A_d^{(0,12)}=\sum_{h=0}^{12}a_{dh}$和16时前实际处理量$P_d^{(0,15)}=\sum_{h=0}^{15}p_{dh}$，程序验证$P_d^{(0,15)}-A_d^{(0,12)}\ge0$。该差值越大，说明排班对早间波动越有缓冲。附件结果同时保存了每个小时的$p_{dh}$、$b_{dh}$和容量，可追溯到任意一天和任意小时。

\begin{figure}[htbp]
\centering
\includegraphics[width=.88\textwidth]{figures/constraint_increment.pdf}
\caption{问题二约束相对于问题一产生的逐日人数增量}
\label{fig:increment}
\end{figure}
